\section{Skill and flexibility are independent axes}
\label{sec:axes}

\begin{table}[t]
\centering
\small
\begin{tabular}{lrrl}
\toprule
Correlation against forecast skill & $r$ & $n$ & leave-one-out range \\
\midrule
\emph{Panel A --- 18 BDG2 buildings, three countries} &  &  &  \\
\quad cooling-degree-day share & +0.010 & 18 & -0.308 to +0.077 \\
\quad controllable (HVAC) fraction & +0.283 & 18 & +0.029 to +0.378 \\
\quad mean outdoor temperature & -0.012 & 18 & -0.266 to +0.051 \\
\quad median load & -0.014 & 18 & -0.189 to +0.120 \\
\midrule
\emph{Panel B --- 6 Chinese provinces, independent replication} &  &  &  \\
\quad cooling-degree-day share & +0.054 & 6 & -0.377 to +0.582 \\
\quad mean outdoor temperature & +0.216 & 6 & -0.290 to +0.782 \\
\quad load--temperature correlation & -0.060 & 6 & -0.619 to +0.713 \\
\midrule
\emph{controllable fraction} vs climate (not vs skill) & \textbf{+0.489} & 18 & -- \\
\bottomrule
\end{tabular}
\caption{Forecast skill is uncorrelated with climate, controllable fraction, temperature and load size, and every leave-one-out range straddles or nearly straddles zero --- which at these sample sizes is what a null looks like. Panel B is the replication on a panel the null was not fitted to. The final row is the one relationship in the study that is real, and note what it is between: available flexibility tracks climate strongly, while predictability does not track anything. Skill and flexibility are independent deployment axes.}
\label{tab:correlations}
\end{table}


\paperfig{fig_null.pdf}{(a) The null. (b) The same null on a panel it was not fitted to, in a
different country, year and data-generating process; panels (a) and (b) share a
vertical scale. (c) The same eighteen buildings, the same horizontal axis, and
the one relationship in the study that is real. How much load there is to shift
tracks climate. How well it can be predicted does not.}{fig:null}

We expected forecast skill --- measured with pinball loss, a strictly proper
scoring rule \citep{gneiting2007strictly}, against each series' own
seasonal-naive baseline --- to increase with cooling load, on the reasoning that
weather-driven demand is the learnable part. Table~\ref{tab:correlations} and
Figure~\ref{fig:null} report the opposite: a null. Every correlation against
skill is indistinguishable from zero and every leave-one-out range straddles or
nearly straddles the sign line, which at these sample sizes is what a null
result looks like. It replicates at $r=+0.054$ on the Chinese provincial panel,
which was added afterwards and spans a wider climate range than the whole
European set.

The controllable fraction, by contrast, tracks climate strongly ($r=+0.489$),
running from $37.6\%$ of the meter in Phoenix to $0.8\%$ in London.

The deployment consequence is a screening rule. Do not ask whether a site's load
is predictable; ask how much of it is weather-driven. A Dublin office is highly
predictable (skill $+0.513$) with roughly $3\%$ controllable load, and there is
nothing to sell there. Conflating the axes inverts the decision. For the target
market this is favourable: Delhi demand correlates $+0.59$ with outdoor
temperature against $-0.64$ for France, placing India at the useful end of the
axis that matters.
